\section{Introduction}
\label{sec:intro}

On-call engineers do not lack alerts; they lack \emph{context}. When a pager fires at 3~AM, the site reliability engineer (SRE) knows the system is unhealthy --- modern observability stacks excel at telling them \emph{that}. What the SRE actually needs, in the first minute of an incident, is the answer to a different question: \emph{``has this happened before, and what did we do about it?''} Jira issue trackers contain rich, structured knowledge about past incidents and their resolutions, but they are unindexed against live telemetry. Mature observability stacks bridge metrics to dashboards, but not metrics to \emph{past tickets}.

This indexing gap shows up as wasted engineer-minutes in every post-incident review. The on-call engineer spends the first 5--30 minutes of an incident manually correlating dashboard signals against half-remembered past pages, even when a directly applicable Jira ticket exists. The bottleneck is no longer detection --- it is \emph{diagnosis-by-citation}: finding the past ticket that explains the current alert.

We pose retrieval-augmented incident triage as a complementary capability to anomaly detection. Anomaly detection answers ``is something wrong?'' Retrieval-augmented triage answers ``what is wrong, and have we seen this before?'' The latter directly reduces the most expensive step of incident response: the diagnosis phase, where an engineer scans logs, traces, and past tickets trying to recognize a recurring pattern.

We make this concrete by building a deterministic skill-chain pipeline that consumes a 5-minute telemetry window (logs, traces, Kubernetes events, and 94 derived numeric features) and produces a triage decision plus a ranked list of past Jira tickets, with each match explained by shared graph entities (service, component, error class). The retrieval side uses BM25, log-signature similarity, and a cross-encoder reranker over a corpus of 347 engineer-voice Jira tickets. We then characterize this system empirically on a synthetic-but-realistic dataset of 2940 test windows --- and crucially, we measure retrieval quality \emph{as a function of how many compatible prior tickets exist in memory}. This deployment-history axis is, to our knowledge, novel in the AIOps literature.

\subsection{Research questions}

This paper answers six research questions:

\begin{itemize}
    \item \textbf{RQ1.} How does retrieval-augmented diagnosis quality vary with the number of compatible prior tickets accumulated in memory?
    \item \textbf{RQ2.} Does anomaly detection improve with Jira memory, or is it an orthogonal capability?
    \item \textbf{RQ3.} Can fine-tuning a domain-specific cross-encoder reranker on $(\textit{window}, \textit{ticket})$ pairs improve retrieval over an off-the-shelf reranker?
    \item \textbf{RQ4.} Which telemetry channel --- logs, traces, or Kubernetes events --- contributes most to retrieval signal?
    \item \textbf{RQ5.} How does retrieval degrade as memory becomes noisier with irrelevant ``distractor'' tickets?
    \item \textbf{RQ6.} What is the engineer-time-saved when a top-$K$ retrieval head is added to an existing triage system?
\end{itemize}

\subsection{Contributions}

\begin{enumerate}
    \item A retrieval-augmented triage architecture that composes BM25, cross-encoder reranking, and a multi-channel evidence graph over a humanized Jira memory corpus.
    \item An empirical characterization of how retrieval quality scales with deployment history --- the central result of the paper, supported by 95\% paired bootstrap confidence intervals over 1000 resamples.
    \item A methodological correction: under the standard $\text{Recall@}K = |\text{top-}K \cap \text{gold}| / |\text{gold}|$ definition, retrieval quality mechanically drops as $|\text{gold}|$ grows beyond $K$. We argue Hit@$K$ (binary recall at top-$K$) is the right metric for ``did the engineer find a relevant ticket'' and show it produces strictly monotone depth-scaling curves.
    \item A fine-tuned cross-encoder reranker (MS-MARCO MiniLM, 12{,}641 training pairs) and a per-channel ablation showing log-style content carries the dominant retrieval signal in our V2 humanized corpus.
    \item A distractor-robustness sweep showing Hit@5 is exactly invariant across $\{0\%, 10\%, 25\%, 50\%\}$ memory occupancy by plausibly-realistic but irrelevant tickets.
    \item An honest negative result: cold-start novelty detection is broken in the current pipeline (0 of 333 cold-start anomalies correctly flagged), motivating future work on calibrated novelty heads.
    \item An open dataset, pipeline, and reproducibility infrastructure: every reported number is traceable to a specific git commit via an automatic per-pipeline training-run registry that writes \texttt{config.json}, \texttt{metrics.json}, \texttt{predictions.jsonl}, and \texttt{train.log} for each fit.
\end{enumerate}

The remainder of the paper is organized as follows. Section~\ref{sec:related} surveys related work. Section~\ref{sec:system} describes the pipeline architecture and the humanized corpus. Section~\ref{sec:method} describes the dataset, metrics, and statistical protocol --- including the Hit@$K$ correction. Section~\ref{sec:results} presents the empirical results for all six research questions. Sections~\ref{sec:discussion} and~\ref{sec:limitations} discuss interpretations and threats to validity. Section~\ref{sec:conclusion} concludes.
